\begin{abstract}
Centroidal Voronoi tessellation (CVT) via Lloyd iteration produces high-quality isotropic surface remeshing but is computationally expensive: the $K$-nearest-neighbor (KNN) query performed for every site at every iteration accounts for 60--80\% of GPU runtime, and all prior methods recompute it unconditionally regardless of whether a site has converged. We observe that convergence in Lloyd iteration is spatially non-uniform---sites in flat regions stabilize early while sites near high curvature oscillate persistently---and that detecting convergence reliably requires accounting for this non-uniformity. A naive freeze policy that removes converged sites from KNN queries produces a 64\% false-freeze rate at curved regions because displacement alone cannot distinguish genuine convergence from transient pauses in oscillation.

We present two synergistic contributions. First, a curvature-adaptive freeze policy that tests both displacement stability and KNN topology stability, requiring both conditions to hold for a curvature-dependent number of consecutive iterations before freezing a site. This reduces false-freeze rates to below 2\% while progressively freezing 80--98\% of sites depending on mesh scale. Second, a reusable bitonic KNN structure that translates the growing frozen fraction into proportional GPU speedup through frozen-mask compaction, warm-starting from previous results, and periodic refresh to bound stale-KNN error.

We evaluate on 11 meshes spanning 35K to 2.84M vertices. On 8 meshes up to 544K vertices, our method is $3.1$--$10.4\times$ faster than Geogram RVD-CVT---a leading CPU Delaunay-based implementation---when both run 250 Lloyd iterations with matched site counts, while producing comparable element quality ($Q_{\mathrm{avg}}$ 0.926 vs.\ 0.929). On 3 meshes at 1.2--2.8M vertices, the freeze policy achieves $3.7$--$4.2\times$ additional speedup over KNN reuse alone, reducing a 2.84M-vertex mesh from 16.7 minutes to 4.0 minutes. Quality metrics track the unfrozen baseline within 0.1\% across all meshes. To our knowledge, this is the first KNN-based GPU CVT demonstrated at million-vertex scale with practical runtime.
\end{abstract}
